\documentclass[11pt]{article}
\usepackage[margin=1in]{geometry}
\usepackage{amsmath,amssymb,bm}
\usepackage{booktabs}
\usepackage{array}
\usepackage{hyperref}
\usepackage{microtype}
\usepackage{siunitx}

\title{One-Dimensional Fatigue Modeling of Pure Aluminum:\\Deterministic Calibration and a Mechanically Generated Spacing Probability}
\author{Donghyeok Kim \and Minsoo Kang \and Junhyeok Park}
\date{Working manuscript, September 2026}

\newcommand{\E}{\mathrm{E}}
\newcommand{\Var}{\mathrm{Var}}
\newcommand{\Cov}{\mathrm{Cov}}

\begin{document}
\maketitle

\begin{abstract}
We formulate a one-dimensional fatigue-initiation model in which the probability state of inter-layer spacing is generated by deterministic nonlinear mechanics rather than selected from an assumed statistical family or empirical life distribution. A finite generalized Lennard--Jones chain supplies the microscopic trajectories. The empirical phase-space measure and its spacing marginal define the probability state exactly, and an exact transport identity yields reduced fields consisting of spacing density, conditional mean spacing rate, and conditional spacing-rate variance. The probability state therefore depends on the external-stress history and the initial microscopic state through the deterministic flow map. Mean spacing and mean interaction energy follow as direct moments, whereas crack initiation is represented by a kinetic first-passage construction at the loss of local normal tangent stiffness. We further distinguish instantaneous high-spacing tail mass from cumulative first-passage mass. Numerical audits show that the instantaneous tail can grow and subsequently decrease through return crossings, while first-passage mass is monotone by construction; external work and conservative mechanical-energy change agree to numerical accuracy. A registry coordinate is retained only as a non-mainline plasticity extension: normal opening lowers the registry barrier in the reduced multilayer potential, but perfect symmetry does not generate registry motion, and permanent plasticity is not claimed without a physical symmetry-breaking and irreversible mechanism. The present model therefore establishes a mechanically generated probability framework for fatigue-relevant state evolution while leaving laboratory time-scale bridging, irreversible dissipation, specimen-scale initial-state measures, and experimental validation as open tasks.
\end{abstract}

\section{Scope and central claim}
The present work is restricted to a one-dimensional normal-loading baseline. It does not reconstruct an FCC lattice, does not postulate a Weibull, Gaussian, Boltzmann, or other lifetime distribution, and does not introduce a phenomenological scalar fatigue-damage variable. The central modeling statement is
\begin{equation}
\sigma_{\rm ext}(0:t),\;\Gamma_0
\longrightarrow
\{a_i(t)\}
\longrightarrow
P_a(a,t),
\end{equation}
where \(\Gamma_0\) denotes the complete initial microscopic state and \(a_i\) are the evolving local normal spacings. Fatigue initiation is then formulated from this mechanically generated probability state and its conditional velocity statistics.

The current paper separates four observables or goals:
\begin{align}
G1:&\quad \bar a(t),\\
G2:&\quad \bar U(t),\\
G3:&\quad E_{\rm hyst}(t)=\int_0^t \dot D_{\rm irr}(t')\,dt',\qquad \dot D_{\rm irr}\ge 0,\\
G4:&\quad \text{first passage to a mechanically defined normal-instability threshold.}
\end{align}
The conservative normal chain presently closes G1 and G2 and supplies a mechanically consistent G4 construction. A nonzero G3 remains an open physical problem; external work, dynamic non-retracing, and cumulative first passage are not identified with irreversible thermodynamic dissipation.

\section{Normalized generalized-Lennard--Jones normal chain}
\subsection{Scaling}
Let \(a_0\) be the reference normal spacing and define
\begin{equation}
\lambda=\frac{a}{a_0},\qquad \tau=\frac{t}{t_0},
\end{equation}
with
\begin{equation}
t_0=\sqrt{\frac{m_a a_0}{E A_0}}.
\end{equation}
The current stress bridge is a calibration definition,
\begin{equation}
q(\tau)=\frac{F_{\rm ext}}{EA_0}=\frac{\sigma_n(t)}{E},
\end{equation}
not a universal atomistic identity.

The normalized generalized-Lennard--Jones potential is
\begin{equation}
\phi(\lambda)
=\frac{\lambda^{-m}}{m(m-n)}
-\frac{\lambda^{-n}}{n(m-n)},
\end{equation}
so that
\begin{equation}
\phi'(\lambda)
=\frac{\lambda^{-n-1}-\lambda^{-m-1}}{m-n},
\end{equation}
\begin{equation}
\phi''(\lambda)
=\frac{(m+1)\lambda^{-m-2}-(n+1)\lambda^{-n-2}}{m-n},
\end{equation}
with \(\phi'(1)=0\) and \(\phi''(1)=1\). The active numerical calibration uses \(m=12.19\) and \(n=6\).

\subsection{Finite-chain equations}
Let \(x_0=0\) and let \(x_1,\ldots,x_M\) denote normalized node positions. Define
\begin{equation}
\lambda_i=x_i-x_{i-1},\qquad c_i=\dot\lambda_i,
\end{equation}
where the dot denotes differentiation with respect to \(\tau\). The potential is
\begin{equation}
V^*=\sum_{i=1}^{M}\phi(\lambda_i).
\end{equation}
The node equations are
\begin{equation}
\ddot x_j=\phi'(\lambda_{j+1})-\phi'(\lambda_j),\qquad j=1,\ldots,M-1,
\end{equation}
\begin{equation}
\ddot x_M=-\phi'(\lambda_M)+q(\tau).
\end{equation}
The exact mechanical-energy balance is
\begin{equation}
E_{\rm mech}^*=\frac12\sum_{j=1}^{M}\dot x_j^2+\sum_{i=1}^{M}\phi(\lambda_i),
\end{equation}
\begin{equation}
\frac{dE_{\rm mech}^*}{d\tau}=q(\tau)\dot x_M.
\label{eq:energybalance}
\end{equation}
Hence a conservative internal model may gain or lose mechanical energy under external cyclic forcing. Such energy exchange is not, by itself, irreversible dissipation.

\section{Mechanically generated probability state}
\subsection{Finite empirical measure}
For one deterministic chain, the spacing probability measure is defined exactly by
\begin{equation}
P_{\lambda,M}(\lambda,\tau)
=\frac1M\sum_{i=1}^{M}\delta\!\left[\lambda-\lambda_i(\tau)\right].
\label{eq:empiricalP}
\end{equation}
This is an empirical atomic probability measure, not an assumed smooth probability-density family. The physical and dimensionless densities satisfy the Jacobian relation
\begin{equation}
P_a(a,t)\,da=P_\lambda(\lambda,t)\,d\lambda,
\end{equation}
so that
\begin{equation}
P_\lambda(\lambda,t)=a_0P_a(a_0\lambda,t),
\qquad
P_a(a,t)=\frac1{a_0}P_\lambda\!\left(\frac{a}{a_0},t\right).
\end{equation}

For an initial-state measure \(\mu_0\), let the full deterministic flow under the applied history \(q\) be
\begin{equation}
\Gamma(\tau)=\Phi^q_{\tau,\tau_0}(\Gamma_0).
\end{equation}
If \(\Lambda_i(\tau;\Gamma_0,q)\) is the spacing produced by that flow, then the exact push-forward form is
\begin{equation}
P_\lambda(\lambda,\tau\mid q,\mu_0)
=\frac1M\sum_{i=1}^{M}\int
\delta\!\left[\lambda-\Lambda_i(\tau;\Gamma_0,q)\right]\mu_0(d\Gamma_0).
\label{eq:pushforwardP}
\end{equation}
For a single deterministic initial condition, \(\mu_0=\delta_{\Gamma_0^*}\), Eq.~\eqref{eq:pushforwardP} reduces to Eq.~\eqref{eq:empiricalP}. Thus the form of \(P\) is fixed by the initial state, the external-stress history, and the nonlinear mechanics.

\subsection{Phase-space measure and exact transport}
Define the empirical phase-space measure
\begin{equation}
F_M(\lambda,c,\tau)
=\frac1M\sum_i
\delta(\lambda-\lambda_i)\delta(c-c_i).
\end{equation}
Let
\begin{equation}
\mathcal G_M(\lambda,c,\tau)
=\frac1M\sum_i\ddot\lambda_i
\delta(\lambda-\lambda_i)\delta(c-c_i).
\end{equation}
Then, distributionally,
\begin{equation}
\partial_\tau F_M+\partial_\lambda(cF_M)+\partial_c\mathcal G_M=0.
\label{eq:exacttransport}
\end{equation}
For a smoothed representation one may write \(\mathcal G=AF\), where
\begin{equation}
A(\lambda,c,\tau)=
\E[\ddot\lambda_i\mid \lambda_i=\lambda,c_i=c].
\end{equation}
This is an exact projected identity; it is not an autonomous closure assumption.

\section{Reduced moment fields \texorpdfstring{\(P,u,\Theta\)}{P,u,Theta}}
Define
\begin{equation}
u(\lambda,\tau)=\E[c\mid\lambda,\tau],
\end{equation}
\begin{equation}
\Theta(\lambda,\tau)=\Var(c\mid\lambda,\tau),
\end{equation}
and the third central conditional moment
\begin{equation}
C_3(\lambda,\tau)=\E[(c-u)^3\mid\lambda,\tau].
\end{equation}
The zeroth and first moment equations give
\begin{equation}
\partial_\tau P+\partial_\lambda(Pu)=0,
\qquad J=Pu,
\label{eq:continuity}
\end{equation}
\begin{equation}
\partial_\tau(Pu)+\partial_\lambda\left[P(u^2+\Theta)\right]=P\mathcal A,
\end{equation}
where
\begin{equation}
\mathcal A(\lambda,\tau)=\E[\ddot\lambda_i\mid\lambda_i=\lambda,\tau].
\end{equation}
Equivalently,
\begin{equation}
D_\tau u
=\mathcal A-\frac1P\partial_\lambda(P\Theta),
\qquad
D_\tau=\partial_\tau+u\partial_\lambda.
\end{equation}
The variance equation is
\begin{equation}
D_\tau\Theta
+2\Theta\partial_\lambda u
+\frac1P\partial_\lambda(PC_3)
=2\Psi,
\label{eq:theta}
\end{equation}
with
\begin{equation}
\Psi(\lambda,\tau)=\Cov(c,\ddot\lambda\mid\lambda,\tau).
\end{equation}
Therefore neither \(C_3=0\) nor \(\Psi=0\) is adopted without additional justification. The triplet \((P,u,\Theta)\) is a history-bearing reduced descriptor, not an autonomous closed Markov state.

Where \(P>0\) and \(\Theta>0\), the momentum identity can be rearranged to the exact smooth density-shape relation
\begin{equation}
\Theta\,\partial_\lambda\ln P
=\mathcal A-D_\tau u-\partial_\lambda\Theta,
\end{equation}
which integrates to
\begin{equation}
P(\lambda,\tau)
=\frac{\mathcal N_P(\tau)}{\Theta(\lambda,\tau)}
\exp\left[
\int_{\lambda_*}^{\lambda}
\frac{\mathcal A(\eta,\tau)-D_\tau u(\eta,\tau)}{\Theta(\eta,\tau)}\,d\eta
\right].
\label{eq:pshape}
\end{equation}
This equation is a representation of a mechanically generated density, not a chosen statistical family. In the degenerate limit \(\Theta\to0\), the undivided moment equations must be used.

\section{Instantaneous tails and cumulative first passage}
\subsection{Snapshot tail flux}
For a fixed threshold \(\lambda_*\), define the instantaneous high-spacing tail
\begin{equation}
Q_*(\tau)=\int_{\lambda_*}^{\infty}P_\lambda(\lambda,\tau)\,d\lambda.
\end{equation}
If the flux vanishes at the upper boundary, Eq.~\eqref{eq:continuity} gives
\begin{equation}
\frac{dQ_*}{d\tau}=J(\lambda_*,\tau).
\label{eq:tailflux}
\end{equation}
Over cycle \(n\),
\begin{equation}
Q_{*,n+1}-Q_{*,n}
=\int_{\tau_n}^{\tau_{n+1}}J(\lambda_*,\tau)\,d\tau.
\label{eq:cycletail}
\end{equation}
With the phase-space density, the flux can be split into outward and return contributions,
\begin{equation}
J_+=\int_0^\infty cF(\lambda_*,c,\tau)\,dc,
\end{equation}
\begin{equation}
J_-=-\int_{-\infty}^0 cF(\lambda_*,c,\tau)\,dc,
\end{equation}
so that \(\dot Q_*=J_+-J_-\). Consequently, instantaneous tail mass is not a cumulative damage variable and need not be monotone.

A related exact moment identity is
\begin{equation}
\frac{d}{d\tau}\Var(\lambda)
=2\Cov(\lambda,u),
\end{equation}
under vanishing boundary contributions. Distribution broadening therefore follows from a spacing--conditional-velocity covariance rather than from an assumed diffusion coefficient.

\subsection{First-passage memory}
For each spacing, define
\begin{equation}
\tau_i^*(\lambda_*)
=\inf\{\tau\ge\tau_0:\lambda_i(\tau)\ge\lambda_*\}.
\end{equation}
The cumulative crossing fraction is
\begin{equation}
F_*(\tau;\lambda_*)
=\frac1M\sum_i I[\tau_i^*(\lambda_*)\le\tau].
\label{eq:firstpassfrac}
\end{equation}
Unlike \(Q_*\), this first-passage observable is nondecreasing by definition. It records history even if a spacing later returns below the threshold. This mathematical memory does not require assigning positive irreversible dissipation to the conservative chain.

\section{G1, G2, and G4}
The mean physical spacing is
\begin{equation}
\boxed{\bar a(t)=\int aP_a(a,t)\,da},
\end{equation}
which, for the finite empirical measure, equals the arithmetic mean of the microscopic spacings. In dimensionless form,
\begin{equation}
\bar\lambda=\int\lambda P_\lambda\,d\lambda,
\qquad \bar a=a_0\bar\lambda.
\end{equation}

The mean intrinsic normal energy is
\begin{equation}
\boxed{
\bar U(t)=U_{\rm ref}\int[\phi(\lambda)-\phi(1)]P_\lambda(\lambda,t)\,d\lambda
}.
\end{equation}
It follows directly from the same microscopic potential that drives the chain.

For normal crack-initiation first passage, the local tangent stiffness vanishes at
\begin{equation}
\phi''(\lambda_c)=0,
\end{equation}
which gives
\begin{equation}
\boxed{
\lambda_c=\left(\frac{m+1}{n+1}\right)^{1/(m-n)}
}.
\end{equation}
For \(m=12.19\) and \(n=6\),
\begin{equation}
\lambda_c\approx1.1077715386.
\end{equation}
The local first-passage time is
\begin{equation}
\tau_i^c=\inf\{\tau\ge\tau_0:\lambda_i(\tau)\ge\lambda_c\}.
\end{equation}
The instantaneous critical tail
\begin{equation}
Q_c(\tau)=\int_{\lambda_c}^{\infty}P_\lambda(\lambda,\tau)\,d\lambda
\end{equation}
is not the cumulative crack-initiation probability. A kinetic survivor density with an absorbing outflow condition at \(\lambda_c\) supplies the first-passage survival probability and hazard. The specimen-level event is the minimum first-passage time over the correlated microscopic system, not an independent-cell product unless independence is separately justified.

\section{Registry coordinate as a plasticity extension, not the baseline}
The active paper does not require a registry variable. However, a reduced multilayer extension can be retained for studying how normal opening changes the accessibility of registry motion. Its intrinsic energy is
\begin{equation}
U_0(a,s)
=\sum_{k\ge1}\sum_{p\in\mathbb Z}
 v_{m,n}\!\left(\sqrt{k^2a^2+(pb+s)^2}\right).
\label{eq:u0as}
\end{equation}
For a stable registry \(s_0\) and neighboring saddle \(s^\dagger\), define
\begin{equation}
\Delta G_s(a)=U_0(a,s^\dagger)-U_0(a,s_0).
\label{eq:registrybarrier}
\end{equation}
A direct-sum diagnostic using the historical normalized parameters \(m=12\), \(n=6\), \(b=\sigma_{LJ}=\epsilon_{LJ}=1\), and \(s_0/b=1/2\) gives a strong reduction of this barrier with normal opening: \(\Delta G_s\) decreases from approximately \(1.48754\) at \(a=0.99196\) to \(0.13312\) at the normal-curvature loss \(a\approx1.13069\). This is a barrier-lowering result, not proof of plastic flow.

If \(g=\Delta G_s(a)\) is monotone over the range of interest, the mechanically generated spacing density induces a barrier density by push-forward,
\begin{equation}
P_G(g,t)
=\int P_a(a,t)\delta[g-\Delta G_s(a)]\,da,
\end{equation}
so that locally
\begin{equation}
P_G(g,t)
=\frac{P_a(a(g),t)}{|d\Delta G_s/da|_{a(g)}}.
\end{equation}
Thus a high-\(a\) population maps to a low-barrier population when \(d\Delta G_s/da<0\).

Two caveats are essential. First, perfect registry symmetry remains invariant. If \(\partial_sU_0(a,s_0)=0\) for every \(a\) and
\begin{equation}
s(0)=s_0,\qquad \dot s(0)=0,
\end{equation}
then the conservative equation \(\mu_s\ddot s=-\partial_sU_0\) admits the exact solution \(s(t)=s_0\). A physical seed or symmetry-breaking mechanism is required before registry motion can occur. Second, the local registry natural frequency obtained from the present reduced coordinate is on an atomic, approximately THz scale under the retained aluminum calibration; direct principal parametric resonance with laboratory-scale fatigue loading is therefore not adopted as the activation mechanism.

If registry motion is later activated, an unwrapped coordinate may be written
\begin{equation}
s=s_0+zb+\widetilde s,
\end{equation}
with well population
\begin{equation}
p_z(t)=\iint_{\mathcal W_z}P_{as}(a,s,t)\,da\,ds.
\end{equation}
A barrier approach, an actual saddle crossing, and residual plasticity are distinct events. Permanent plasticity would require a residual change in the well population after unloading and relaxation; it is not established merely by transient access to a low-barrier state.

\section{G3: external work is not irreversible dissipation}
For the current conservative chain, Eq.~\eqref{eq:energybalance} accounts for the complete mechanical-energy change. A load/unload mismatch of \(P\), \(u\), or \(\Theta\), a nonzero cycle work increment, or a cumulative first-passage record does not by itself prove
\begin{equation}
\dot D_{\rm irr}>0.
\end{equation}
The present baseline therefore retains
\begin{equation}
\dot D_{\rm irr}=0,
\qquad
E_{\rm hyst}=0
\end{equation}
as its thermodynamic statement. If a future physical irreversible force \(r_j^{\rm irr}\) is derived, the balance would become
\begin{equation}
\frac{dE_{\rm mech}^*}{d\tau}
=q\dot x_M-\dot D_{\rm irr}^*,
\end{equation}
with
\begin{equation}
\dot D_{\rm irr}^*
=-\sum_j r_j^{\rm irr}\dot x_j\ge0.
\end{equation}
No damping coefficient, white-noise bath, or Markov reduction is inserted solely to manufacture G3.

\section{Numerical verification}
\subsection{Density-shape reconstruction}
A 512-atom reconstruction audit of Eq.~\eqref{eq:pshape}, using numerical smoothing only to estimate finite-sample conditional fields, reported a maximum \(L^1\) density error of \(0.076255\), a maximum Kolmogorov--Smirnov distance of \(0.038126\), and a maximum relative mean-energy error of \(3.046\%\) over the tested phases. The Gaussian kernel used in that audit is a numerical estimator rather than a physical probability assumption.

\subsection{Same-force non-retracing}
A separate 512-atom matched-force audit compared loading and unloading states at almost identical applied force. The reduced descriptors were not single-valued functions of instantaneous force. For example, at the approximately \(150\,\mathrm{MPa}\) matched level the density comparison gave \(L^1\approx0.7023\) and KS \(\approx0.2061\), while the normalized differences in \(u\) and \(\Theta\) were approximately \(1.406\) and \(1.828\), respectively. The chain remained conservative: per-cycle external-work and mechanical-energy increments agreed to errors of order \(10^{-13}\)--\(10^{-12}\). The result is therefore dynamic history dependence, not evidence of irreversible G3 dissipation.

\subsection{Tail-flux and first-passage audit}
We additionally audited a six-cycle protocol with 512 atoms, \(\Delta\tau=0.01\), mean and amplitude each mapped from \(100\,\mathrm{MPa}\) through \(\sigma/E\), \(\omega^*=0.02\), and a two-cycle smooth ramp. At the illustrative threshold \(\lambda_*=1.004\), the cycle-end snapshot tail fractions were
\begin{equation}
0,\;0,\;0.037182,\;0.266145,\;0.256360,\;0.146771,
\end{equation}
for cycles 1--6. Hence snapshot tail mass is explicitly non-monotone. The corresponding cumulative first-passage fractions were
\begin{equation}
0,\;0,\;0.043053,\;0.491194,\;0.704501,\;0.743640,
\end{equation}
which are monotone. Direct crossing counting verified Eq.~\eqref{eq:cycletail}: during cycles 3--6 the upward/downward crossing pairs were \((22,3)\), \((251,134)\), \((358,363)\), and \((330,386)\), giving net changes \(19\), \(117\), \(-5\), and \(-56\), exactly equal to the changes in the number of spacings above threshold. The largest absolute mismatch between cumulative external work and mechanical-energy change was approximately \(9.15\times10^{-12}\).

No spacing reached the actual normal-instability threshold \(\lambda_c\) in this six-cycle low-load proof-of-principle protocol. Therefore the audit validates the tail/first-passage bookkeeping and conservative energy accounting, not a calibrated fatigue-life prediction. The dimensionless frequency \(\omega^*=0.02\) is not mapped directly to laboratory fatigue frequency.

\section{What is established and what remains open}
The present formulation establishes the following mathematical chain:
\begin{equation}
\sigma_{\rm ext}(0:t),\Gamma_0
\rightarrow
\text{deterministic normal trajectories}
\rightarrow
P_a(a,t),u_a(a,t),\Theta_a(a,t)
\rightarrow
\{G1,G2,\text{first-passage }G4\}.
\end{equation}
It also establishes that instantaneous tail evolution is governed by probability flux and that cumulative first passage retains threshold-crossing history without requiring an empirical damage variable. The registry extension provides a mechanically defined route from normal spacing to a distribution of registry barriers, but it does not yet supply permanent plasticity.

The principal open problems are:
\begin{enumerate}
\item derive or identify a physically justified microscopic irreversible mechanism if nonzero G3 is required;
\item bridge the atomic dynamic time scale to laboratory fatigue frequencies without assigning an artificial slow registry inertia;
\item construct and calibrate the physically appropriate initial full-state measure \(\mu_0\) and specimen-scale correlation structure;
\item test whether registry barrier lowering actually produces saddle crossing and residual well-index change under a justified symmetry-breaking seed;
\item validate the predicted spacing distributions, tail fluxes, and first-passage trends experimentally.
\end{enumerate}

\section{Reproducibility paths}
The tail/first-passage audit is generated by
\texttt{simulations/verify\_probability\_tail\_first\_passage.py}
on the \texttt{numerical-fem} branch. Its tabulated output is stored under
\texttt{results/data/probability\_tail\_first\_passage/}, with an interpretation report at
\texttt{results/reports/PROBABILITY\_TAIL\_FIRST\_PASSAGE.md}.
The density-shape and same-force history audits are stored under
\texttt{results/data/theta\_density\_reconstruction/} and
\texttt{results/data/theta\_same\_force\_history/}, respectively.

\section{Conclusion}
A probability density for local normal spacing need not be assumed in advance. In the present one-dimensional model it is the push-forward of deterministic nonlinear chain mechanics under a specified initial state and external-stress history. This construction gives exact probability transport, reduced conditional moment equations, direct mean-spacing and energy observables, and a kinetic first-passage definition of normal instability. Numerical audits support the distinction between reversible snapshot-tail evolution and monotone first-passage memory while preserving conservative energy balance. A registry extension shows that normal opening can strongly reduce a slip barrier, but symmetry breaking, residual plasticity, laboratory time-scale mapping, irreversible dissipation, and experiment remain unresolved. These boundaries define the next validation stage rather than additional assumed constitutive laws.

\end{document}
